\section{Introduction}\label{sec:introduction}

The Penrose inequality relates the total mass of an asymptotically flat initial-data set to the area required to enclose its black-hole boundary. Penrose proposed this relation as a consequence of gravitational collapse and cosmic censorship \cite{Penrose1973}. In three dimensions, Huisken and Ilmanen proved the Riemannian inequality for a connected outermost minimal component by inverse mean curvature flow, and Bray treated total horizon area with multiple components by conformal deformation \cite{HuiskenIlmanen2001,Bray2001}. Bray and Lee extended the numerical Riemannian inequality to spatial dimensions below eight \cite{BrayLee2009}. The spacetime problem retains the second fundamental form and the invariant ADM mass, and the charged problem must also account for the conserved electric flux. These distinctions affect both the numerical inequality and the equality statement.

This article proves a reduction from the neutral spacetime inequality to the purely electric inequality in every spatial dimension $n\ge4$. The neutral input is stated in full in Theorem~\ref{thm:neutral-input}. It is used only at its stated stronger-decay, future-trapped, positive-area, future-timelike scope. We use this companion result without independently reproving the unrestricted neutral spacetime Penrose conjecture. Using this input, Theorems~\ref{thm:numerical} and~\ref{thm:rigidity} establish the numerical bound and connected nonextremal rigidity statement formulated below.

The numerical theorem permits disconnected boundary, with a future or past trapping sign chosen independently on each component. It first establishes that the ADM vector is future timelike; write $m$ for its invariant mass. The area is an infimum over \emph{full enclosing cuts}. Every component of a cut, including its contact with the given boundary, is counted, and no smooth minimizer is assumed. If $Q$ is the signed charge and $X_A$ is the $(n-2)$nd power of the corresponding area radius, the sharp assertion is
\begin{equation}\label{intro:bound}
 m\ge |Q|,\qquad X_A\le m+\sqrt{m^2-Q^2}.
\end{equation}
On the outer branch $X_A>|Q|$, the second inequality is equivalent to $m\ge (X_A+Q^2/X_A)/2$; for $0<X_A\le|Q|$, the stated conclusion reduces to the mass-charge bound. This distinction is necessary when the enclosing boundary has several components. In spatial dimension three, Khuri, Weinstein and Yamada proved the corresponding upper area-radius bound for strongly asymptotically flat Riemannian data with possibly disconnected outermost minimal boundary, source-free electric and magnetic fields, and the charged scalar-curvature bound \cite[Theorem~1.1 and Corollary~1.2]{KhuriWeinsteinYamada2017}. The present reduction treats nonzero second fundamental form and spatial dimensions $n\ge4$.

For a connected nonextremal future horizon satisfying the outermost and area-minimizing conditions below, equality determines all the original fields as a global slice of a Reissner--Nordstr\"om--Tangherlini exterior. The argument retains the singular minimizing frontiers that can occur from spatial dimension eight onward, and the equality analysis takes place on the original data throughout.

The proof has four main ingredients. First, a charge-preserving end preparation reduces weakly asymptotically flat data with arbitrary timelike ADM vector to strict data with a static end. A filled scalar system then produces a metric which dominates the original metric and has charged scalar-curvature control on a suitable height region. Its local regularity rests on a special rank-one elliptic structure: after freezing the gradient direction, the derivative in that direction satisfies a divergence-form equation even when the remaining scalar coefficient is merely measurable. This yields the required gradient estimates in every dimension.

Second, a bounded divergence-flux equation converts electric charge into a controlled decrement of ADM energy. Its conformal factor has a lower barrier, so the loss of enclosing area is explicit. A compact pure-trace second fundamental form makes a smooth regular height cut strictly future trapped. Applying the exact neutral input and optimizing one parameter proves \eqref{intro:bound}. The construction permits sources in the artificial filling; the curvature estimate is used only on the original source-free exterior.

Third, equality is studied by varying the original data. A measure multiplier for the full area envelope produces a lapse and shift, together with a positive normal-normal measure in the lapse Hessian. Atomic minimizing sheets are eliminated by a jump identity. Non-atomic sheets carry positive Jacobi fields; a capacity and tangent-cone argument controls their singular ends, and a one-dimensional BV argument removes diffuse multiplier mass on zero-lapse sheets. This is the part of the proof which retains the singular sets allowed in high dimension.

Finally, closed-form electromagnetic variations produce a global electric potential. Conservation of volume along electric flow lines forces alignment of the electric field and shift. A cutoff argument proves staticity across possible interior zeros of the Killing norm; those zeros are then excluded. A transformation to vacuum static data allows conformal doubling and a direct treatment of the compactified point. The resulting global base recovers the original spacelike embedding, its future normal, and the signed normalized electric field. Classical static uniqueness arguments motivate this final stage \cite{BuntingMasoodulAlam1987,GibbonsIdaShiromizu2002}. Higher-dimensional static charged uniqueness was developed by Gibbons, Ida and Shiromizu \cite{GibbonsIdaShiromizu2002ChargedPRD}; the present equality argument must first obtain staticity and then recover the original slice.

The neutral companion paper \cite{NeutralV6S74} and the three-dimensional charged companion paper \cite{ChargedV6S818} supply methodological precedents. Only the displayed numerical neutral theorem is taken as a premise. We reproduce and justify the auxiliary constructions needed for the present dimension, field degree, regularity, and boundary conventions.
